\scanpage{211}
\uline{\unsure{Question|Questions}}
% slanted
\margin{\tikz[baseline=(a.base)]{\draw[gray] (.1,-1.9) .. controls (.25,-.4) and (.17,1.2) .. (.15,2.4); \node[anchor=east,rotate=35](a) at (0,0){À \unsure{revoir|prévoir|voir}};}}

\diagram{\[
\left\{\begin{array}{@{}l@{}}
\left\{\begin{array}{@{}l@{}}
\parbox{.79\linewidth}{pts \inserted{de $X$} à val. dans $k$ \inserted{de car $p>0$} : correspondent aux pt motiviques \inserted{dans $W(k)$ (ou corps des fractions ?)} avec \uline{structure supplémentaire} \tikz[baseline=(a.base)]\node[ellipse,draw,inner sep=1pt](a){$F\deleted{\unsure{F|S}},V$}; [foncteurs des groupes formels]}\\
\begin{tikzcd}[ampersand replacement=\&,column sep=.8em,row sep=.4em]
{} \arrow[dr,squiggly] \& {}\\
{} \& \parbox{.60\linewidth}{pts motiviques généralisés (au \unsure{sens|une} cat. dérivées) dans $k$, avec structure filtrée !}
\end{tikzcd}
\end{array}\right.\\[.6em]
\parbox{.82\linewidth}{pts \inserted{de $X$} à val. dans $k$ de car $0$ : correspondent aux pts motiviques dans $k$, avec \uline{structure filtrée}}
\end{array}\right.
\]}

\diagram{\[
\left|\begin{array}{l}
\text{pts \inserted{de $X$} à valeurs dans $\mathbb C$ mod conj. complexe}\\
\text{(= \unsure{places|pl\textsuperscript{ces}|plus})}\ \longleftrightarrow\
 \text{pts \inserted{motiviques} à valeurs dans }
 \tikz[baseline=(a.base)]\node[draw,ellipse](a){$\mathbb R$};\text{, ou}\\
\quad\text{à valeurs dans }\tikz[baseline=(a.base)]\node[draw,ellipse](a){$\mathbb C$ avec \tikz[baseline=(b.base)]\node[draw,ellipse,inner sep=1pt](b){$f_{\infty}$};};\text{, ou}\\
\quad\text{à valeurs dans }\mathbb Z
\end{array}\right.
\]}

\uline{\deleted{\unsure{Sch de car $p$|sur $X$ de car $p$|de $X$ de car $p$|Sch de cas $p$}}}

\uline{Pts banaux}\qquad Pts à valeurs dans \tikz[baseline=(a.base)]\node[draw,ellipse](a){$\mathbb Z_{\unsure{l|\ell}}$};, $\ell$ premier aux \unsure{car.} résiduelles. ou \tikz[baseline=(a.base)]\node[draw,ellipse](a){$\mathbb Q_{\ell}$};

! \uline{Cas particulier important}\qquad \uline{Pts à valeurs dans corps} fini $\mathbb F_q$ correspondent aux p\textsuperscript{ts} à valeurs dans les \tikz[baseline=(a.base)]\node[draw,ellipse](a){$W(\mathbb F_q)$};, avec structure supplémentaire $(F,V)$
% archive sheet number: 210

